\documentclass[11pt]{article}
\usepackage[margin=1in]{geometry}
\usepackage[T1]{fontenc}
\usepackage{lmodern,microtype,hyperref,enumitem,amsmath,amssymb,xcolor,booktabs}
\usepackage[numbers,sort&compress]{natbib}
\usepackage{url}
\hypersetup{colorlinks=true,urlcolor=blue,citecolor=blue,linkcolor=blue}

\title{Q03 --- Deeper Projects\\\large Is reasoning real, a property of language, a reflection, or an illusion?}
\author{DLS 2026 --- AI Internals Module}
\date{}

\begin{document}
\maketitle

\section*{Framing the Question}

When a large language model emits a chain of intermediate tokens before an
answer, we are tempted to call this ``reasoning.'' But the philosophical
content of that word is contested. Four hypotheses structure this module:
(i) \emph{reasoning is real}---there is a causal, generalizable computation
that mediates premises and conclusions; (ii) \emph{reasoning is a property
of language}---the surface form of trained text induces inference-like
patterns without a separable mechanism; (iii) \emph{reasoning is a
reflection}---the model mirrors statistical regularities of human-produced
argument without executing the argument; (iv) \emph{reasoning is an
illusion}---chain-of-thought (CoT) text is post-hoc rationalization that
does not reflect the model's actual decision process. These are not mutually
exclusive across tasks, and the empirical literature supports a nuanced
view: CoT improves performance on some tasks \citep{wei2022cot,
nye2021scratchpad}, yet verbalized explanations can be unfaithful to the
underlying computation \citep{turpin2023language, lanham2023measuring}, and
apparent reasoning degrades under distribution shift and token-probability
sensitivity \citep{mccoy2023embers}. The projects below ask you to generate
\emph{evidence}, not opinions.

\section{Project A --- Faithfulness Auditing of Chain-of-Thought}

\paragraph{Goal.} Determine whether verbalized reasoning steps are causally
responsible for a model's answer, or whether they are decorative.

\paragraph{Method.} Reproduce and extend the biasing and perturbation
protocols of \citet{turpin2023language} and \citet{lanham2023measuring}.
Construct matched prompt pairs in which a spurious cue (e.g., answer order,
few-shot label bias) is silently inserted. Measure whether the model's
answer flips while its stated rationale fails to mention the cue. Then run
causal interventions: truncate CoT, paraphrase it, inject errors into
intermediate steps, and measure answer sensitivity.

\paragraph{Deliverables.}
\begin{itemize}[leftmargin=1.4em]
  \item A faithfulness benchmark of $\geq 400$ items across $\geq 3$ task
    families (arithmetic, multi-hop QA, logical entailment), with per-item
    bias labels.
  \item A quantitative report giving: answer-flip rate under biasing,
    rationale-mention rate, and an \emph{early-answering} / \emph{added-mistake}
    sensitivity curve per \citet{lanham2023measuring}.
  \item A written argument (3--4 pp.) assigning each task family to one of
    the four hypotheses, with confidence intervals.
\end{itemize}

\section{Project B --- Mechanistic Localization of an Inference Circuit}

\paragraph{Goal.} Test whether a specific reasoning behavior corresponds to
an identifiable, reusable internal computation (the ``reasoning is real''
hypothesis at the mechanism level).

\paragraph{Method.} Choose a narrow, well-posed task (e.g., multi-step
in-context induction, or a bracketed arithmetic subtask). Apply activation
patching and path patching following \citet{wang2023ioi} and the causal
mediation framework of \citet{meng2022rome}. Use attention-pattern analysis
and logit attribution to identify candidate circuits, then validate by
ablation and transfer to held-out instances. Document whether the circuit
generalizes or is template-bound (evidence for ``property of language'').

\paragraph{Deliverables.}
\begin{itemize}[leftmargin=1.4em]
  \item A circuit diagram with the heads/MLPs implicated, plus patching
    effect sizes.
  \item A transfer table: circuit intact across $\geq 4$ structural variants
    (vocabulary swap, length scaling, distractors).
  \item A short claim with falsification criterion: state what observation
    would refute ``this circuit implements the inference.''
\end{itemize}

\section{Project C --- Robustness and the ``Embers'' of Autoregression}

\paragraph{Goal.} Quantify how much apparent reasoning is governed by
token-level probability rather than task logic.

\paragraph{Method.} Following \citet{mccoy2023embers}, construct task
variants that hold logical structure fixed while varying output
probability (common vs.\ rare target strings, common vs.\ rare task
framings). Measure accuracy as a function of estimated output rarity and
task frequency. Contrast with symbolic baselines. Optionally cross-reference
counting/length-sensitivity failures reported in compositional-generalization
work \citep{lake2018scan}.

\paragraph{Deliverables.}
\begin{itemize}[leftmargin=1.4em]
  \item Accuracy-vs-rarity curves with regression coefficients isolating
    the probability effect from the logic effect.
  \item A ``reflection index'': the fraction of variance in correctness
    explained by output probability versus task structure.
  \item Interpretation mapping results onto the four hypotheses.
\end{itemize}

\section{Project D --- Operationalizing the Definition}

\paragraph{Goal.} Rather than settle the metaphysics, propose a
\emph{measurable} definition of reasoning and defend it.

\paragraph{Method.} Draft operational criteria (e.g., causal dependence of
answer on verbalized steps, out-of-distribution generalization, circuit
reusability, invariance to superficial rephrasing). Relate your criteria to
emergent-capability debates \citep{wei2022emergent, schaeffer2023mirage} and
to the self-consistency improvements of \citet{wang2023selfconsistency}.
Apply your criteria to results from Projects A--C.

\paragraph{Deliverables.}
\begin{itemize}[leftmargin=1.4em]
  \item A 2-page definitional proposal with explicit metrics and thresholds.
  \item A scoring of $\geq 2$ models against your criteria, with a verdict
    per hypothesis (real / linguistic / reflective / illusory), per task.
  \item A critical appendix noting where your own definition is itself at
    risk of measuring an illusion.
\end{itemize}

\section*{Assessment}

Projects are graded on experimental validity, honesty about negative
results, and the tightness of the link between \emph{evidence} and
\emph{claim}. A project that convincingly demonstrates ``illusion'' on a
task earns the same credit as one demonstrating ``real'' reasoning;
intellectual direction is not rewarded, only rigor.

\bibliographystyle{plainnat}
\begin{thebibliography}{99}

\bibitem{wei2022cot}
J.~Wei, X.~Wang, D.~Schuurmans, M.~Bosma, B.~Ichter, F.~Xia, E.~Chi,
Q.~Le, and D.~Zhou.
\newblock Chain-of-thought prompting elicits reasoning in large language models.
\newblock In \emph{Advances in Neural Information Processing Systems (NeurIPS)}, 2022.

\bibitem{nye2021scratchpad}
M.~Nye, A.~J. Andreassen, G.~Gur-Ari, H.~Michalewski, J.~Austin, D.~Bieber,
D.~Dohan, A.~Lewkowycz, M.~Bosma, D.~Luan, C.~Sutton, and A.~Odena.
\newblock Show your work: Scratchpads for intermediate computation with language models.
\newblock \emph{arXiv preprint arXiv:2112.00114}, 2021.

\bibitem{turpin2023language}
M.~Turpin, J.~Michael, E.~Perez, and S.~R. Bowman.
\newblock Language models don't always say what they think: Unfaithful
explanations in chain-of-thought prompting.
\newblock In \emph{Advances in Neural Information Processing Systems (NeurIPS)}, 2023.

\bibitem{lanham2023measuring}
T.~Lanham, A.~Chen, A.~Radhakrishnan, B.~Steiner, C.~Denison, D.~Hernandez,
D.~Li, E.~Durmus, E.~Hubinger, J.~Kernion, et~al.
\newblock Measuring faithfulness in chain-of-thought reasoning.
\newblock \emph{arXiv preprint arXiv:2307.13702}, 2023.

\bibitem{mccoy2023embers}
R.~T. McCoy, S.~Yao, D.~Friedman, M.~Hardy, and T.~L. Griffiths.
\newblock Embers of autoregression: Understanding large language models
through the problem they are trained to solve.
\newblock \emph{arXiv preprint arXiv:2309.13638}, 2023.

\bibitem{wang2023ioi}
K.~Wang, A.~Variengien, A.~Conmy, B.~Shlegeris, and J.~Steinhardt.
\newblock Interpretability in the wild: A circuit for indirect object
identification in GPT-2 small.
\newblock In \emph{International Conference on Learning Representations (ICLR)}, 2023.

\bibitem{meng2022rome}
K.~Meng, D.~Bau, A.~Andonian, and Y.~Belinkov.
\newblock Locating and editing factual associations in GPT.
\newblock In \emph{Advances in Neural Information Processing Systems (NeurIPS)}, 2022.

\bibitem{lake2018scan}
B.~M. Lake and M.~Baroni.
\newblock Generalization without systematicity: On the compositional skills
of sequence-to-sequence recurrent networks.
\newblock In \emph{International Conference on Machine Learning (ICML)}, 2018.

\bibitem{wei2022emergent}
J.~Wei, Y.~Tay, R.~Bommasani, C.~Raffel, B.~Zoph, S.~Borgeaud, D.~Yogatama,
M.~Bosma, D.~Zhou, D.~Metzler, et~al.
\newblock Emergent abilities of large language models.
\newblock \emph{Transactions on Machine Learning Research (TMLR)}, 2022.

\bibitem{schaeffer2023mirage}
R.~Schaeffer, B.~Miranda, and S.~Koyejo.
\newblock Are emergent abilities of large language models a mirage?
\newblock In \emph{Advances in Neural Information Processing Systems (NeurIPS)}, 2023.

\bibitem{wang2023selfconsistency}
X.~Wang, J.~Wei, D.~Schuurmans, Q.~Le, E.~Chi, S.~Narang, A.~Chowdhery, and D.~Zhou.
\newblock Self-consistency improves chain of thought reasoning in language models.
\newblock In \emph{International Conference on Learning Representations (ICLR)}, 2023.

\end{thebibliography}

\end{document}